\section{A complete fallback for exceptional compressed-braid factors}
\label{sec:exceptional-cube}

The preceding section left a precise missing implementation step: completing
wider factors after singleton cuts without expanding the already solved
three-braid factors. We now connect a bounded reduced Khovanov cube, with
source-bound linear-algebra proofs, to the native grammar interface. Given
sufficient allowances this makes recognition complete for every supplied braid
grammar. Its running time is polynomial in grammar size plus an exponential
term for each exceptional factor's minority-sign count. The latter parameters
remain uncontrolled for general knots.

\subsection{Dispatch and the source contract}

The existing component check and exact singleton cuts run first. One- and
two-strand factors use exponent sums; three-strand factors use compressed
free-product normal forms. A wider factor with $|e|>m-1$ is already nontrivial.
For every remaining wider factor the host compares its exact binary expanded
length with \texttt{fallback\_max\_crossings}, default 12. Only a factor that
passes this comparison is expanded. All larger residuals remain explicitly
inconclusive; no iteration up to their represented lengths occurs.

Expansion walks the projected grammar and skips zero-length children. This
last detail is necessary: a depth-$s$ grammar of concatenated empty strings
can have $2^s$ occurrences but length zero. With zero children skipped, every
visited node lies on a path to one of the $n$ actual letters, giving an
$O(s(n+1))$ traversal bound. The expansion is checked against its previously
computed length. Three-braid factors never enter this routine, including
factors with exponentially large minority-sign counts.

A forest proof containing an exceptional cube uses version two. Its original
input, singleton cuts, strand intervals and projected grammars are checked
exactly as before. The child proof additionally records the expanded signed
word and strand count, which replay recomputes from that projected grammar.
Thus neither an independently chosen small braid nor a convenient PD diagram
can substitute for the actual factor. The cube operates directly on the braid;
a separate diagram-to-braid conversion is not part of its trusted input.
Version-one proofs remain replayable with their original unsupported-factor
semantics. Changing a version-two cube proof to version one is rejected.

\subsection{The exact reduced cube}

We adapt the explicit braid-resolution construction from report 34's MIT-0
oracle, retaining it as a separately loaded comparison implementation.
For an $m$-strand word of length $n$, vertices comprise the $m$ top endpoints
and two new bottom endpoints per crossing. In each state
$\epsilon\in\{0,1\}^n$, smoothing joins either the two horizontal pairs or
the two vertical pairs. The choice of state bit is reversed for a negative
Artin generator. Joining final endpoints to their initial strand positions
closes the resolution. A disjoint-set computation returns its circles.
The circle containing vertex zero is the marked circle in every state.

Assign each circle a label in $\{1,x\}$ with $x^2=0$, and require the marked
circle's label to be $x$. A state with $c$ circles therefore contributes
$2^{c-1}$ basis elements to cube degree $|\epsilon|$. A cube edge changes
one zero state bit to one and applies the ordinary Frobenius maps
\[
 \begin{array}{ll}
 \mu(1,1)=1,&\mu(1,x)=\mu(x,1)=x,\quad\mu(x,x)=0,\\
 \Delta(1)=1\otimes x+x\otimes1,&\Delta(x)=x\otimes x.
 \end{array}
\]
Unchanged circles keep their labels. The marked-$x$ subcomplex is preserved.
Over $\mathbb F_2$ the cube signs disappear. Homological normalization shifts
and quantum gradings are unnecessary for the total rank used here.

Let $D_h$ be the dimension in raw cube degree $h$, and $r_h$ the rank of its
differential to degree $h+1$. The implementation constructs each differential
as integer bit-vector columns and checks $d^2=0$ before using
\[
 b_h=D_h-r_h-r_{h-1},\qquad
 \dim\widetilde{Kh}(K;\mathbb F_2)=\sum_h b_h.
\]
The same chain check runs during replay. Negative dimensions or zero total
rank on a knot are implementation errors, not verdicts. The unknot test is
total reduced rank one. Kronheimer--Mrowka gives the rational rank-one
criterion~\cite{km}; the universal coefficient theorem bounds rational rank
by mod-two dimension, while the unknot has a single torsion-free reduced
generator. These facts give the same detection criterion over $\mathbb F_2$.
They do not identify the Jones polynomial with a complete decision procedure.

\subsection{A rank proof with no pivot search in replay}

For each matrix column, discovery records a list of earlier basis-column
indices whose XORs eliminate its leading bits, followed by the leading row
of the residual, or $-1$ for zero. The proof stores indices, not copies of all
intermediate bit vectors. Discovery uses a map from leading row to its basis
column to select eliminations.

Replay reconstructs the differential from the source word and reads the
proposed XOR list. Each reference must name a distinct preceding basis
column. A zero residual must have pivot $-1$. A nonzero residual must have
exactly the stated highest set bit, and that row must differ from all earlier
basis pivots. Replay stores this residual as the new basis column. It does
not call the rank producer or select a pivot to eliminate.

The rank assertion follows by two inductions. First, distinct highest set
bits make the retained residual vectors linearly independent: among vectors
in a putative dependence, the largest leading row cannot cancel. Second,
each original column is the XOR of its residual and earlier basis vectors,
while each new residual is in the span of original columns already processed.
The two spans are consequently equal after every column. The number of
retained residuals is the exact rank, including for zero and rectangular
matrices. Both lower and upper rank bounds are proved; a list of independent
columns alone would prove only the lower bound and would be insufficient for
an unknot verdict.

The rank verifier is structurally independent of elimination discovery.
It shares the explicit Frobenius complex construction, expansion helper,
integer arithmetic and $d^2$ check with the producer. This trusted boundary is
intentional and visible. Comparisons with the maintained planar-diagram
scanner test a different geometric construction and reduction algorithm;
comparisons with report 34 additionally check that the adaptation preserves
its original ranks.

\subsection{Budgets across discovery and replay}

The public API adds \texttt{fallback\_max\_crossings=12} and
\texttt{fallback\_max\_generators=200000}, also available as hyphenated CLI
options. The latter is a cumulative allowance across \emph{all} cube
constructions, including reconstruction during replay. A resolution reserves
its $2^{c-1}$ generators before allocating the differential matrices.
The ordinary work callback is polled during expansion, resolution building,
Frobenius edge assembly, the chain check, elimination and rank replay.
String arenas, metadata, JSON transport and these cube operations retain the
same shared work and time budgets.

A factor above the crossing threshold has not started a cube and may remain
as a verified inconclusive residual in a forest proof. Exhaustion of shared
work, generator, node, time or proof-size allowances returns inconclusive
without a partial proof. External cancellation propagates. Setting the
crossing threshold to zero restores the earlier unsupported-factor behavior;
raising it permits a larger exponential computation. Proof byte limits are
checked after construction, so they are not hard peak-memory guarantees.
Neither a generator count nor a work counter claims to count individual
machine instructions or bits allocated by Python.

\subsection{The exceptional-factor theorem is now implemented}

Let $g$ be the original grammar bit size and let $I$ index only wider factors
not decided by the finite checks. For factor $i$, write $n_i,m_i,\kappa_i$ for
its expanded length, strand count and minority-sign count. The previous
singleton and Bennequin arguments give
\[
 n_i\le4\kappa_i,\qquad m_i\le2\kappa_i+1.
\]
The cube has $2^{n_i}$ states. Each resolution has at most $m_i+2n_i$ circles,
since that is its number of vertices, so a deliberately coarse bound for its
total reduced basis size is
\[
 D_i\le 2^{m_i+3n_i-1}=2^{O(\kappa_i)}.
\]
Building all columns and checking $d^2$ use polynomially many bit operations
in $D_i+n_i+m_i$. Elimination performs at most $D_i$ XORs per column on at
most $D_i$ bits. A trace has at most $O(D_i^2)$ indices of $O(\log(D_i+2))$
bits, and replay has the same exponential upper bound. The polynomial costs
of source validation, projections and solved compressed factors are paid
separately. Products of a source-polynomial term and an exponential factor
can be absorbed into an additive bound using $ab\le a^2+b^2$, with only a
constant change in the exponential coefficient. With sufficiently large
allowances, implemented discovery,
certificate production and replay therefore satisfy
\[
 \operatorname{poly}(g)+\sum_{i\in I}2^{O(\kappa_i)}.
\]
All factors eventually decide as the finite allowances increase. The number
of factors is polynomial in $g$. Thus the implementation has a restricted
quasi-polynomial guarantee when
$\max_{i\in I}\kappa_i=O((\log g)^2)$, and a polynomial guarantee when those
counts are $O(\log g)$. There is no restriction on the expanded minority
counts of the solved three-braid factors. This is a complete algorithm for
supplied braid grammars, but an exponential worst case remains. We have not
proved the exceptional-parameter condition or a favorable grammar-construction
bound for arbitrary knot diagrams.

\subsection{Validation and measured scope}

Twelve new tests exercise both knot verdicts, exact source binding, Boolean
field substitutions, altered dimensions and ranks, invalid column spans,
replay without the rank producer, and shared budgets through multiple factors
and replay. All 512 binary $3\times3$ matrices are checked against their
explicitly enumerated column spans. Random braid cubes are compared with the
maintained planar-diagram scanner. A mixed seven-strand forest contains a
$2^{256}$ conjugation sleeve and a nine-crossing exceptional factor; only the
latter is expanded. A separate binary-sized wider factor is declined before
expansion. The empty-subgrammar test ensures that represented emptiness does
not cause exponential traversal.

The complete maintained suite passes 897 tests with Regina in 226.425 seconds.
A separate seeded audit constructs 160 knot closures on three through six
strands with at most ten crossings. All reduced ranks agree with both the
unchanged report-34 oracle and the maintained planar-diagram scanner, with
$d^2$ checked in all three computations. Every public result also replays.
There are 108 unknots and 52 nontrivial knots; 17 public calls require the new
cube fallback, while the other cases are decided by the earlier exact stages.
The audit takes 44.679 seconds and retains every signed input word, rank,
verdict and source hash in \path{data/exceptional-cube-audit.json}.

\begin{center}
\begin{tabular}{@{}lrrrr@{}}
\toprule
Input & whole ms & hybrid ms & paired ratio & generators \\
\midrule
four\_unknot & 159.49 & 161.83 & 1.012 & 4638 \\
four\_knot & 139.22 & 137.07 & 1.016 & 3774 \\
small\_sum\_unknot & 7702.33 & 169.43 & 44.278 & 4638 \\
small\_sum\_knot & 5557.81 & 139.17 & 40.209 & 3774 \\
huge\_sum\_8 & -- & 168.05 & -- & 4638 \\
huge\_sum\_64 & -- & 187.69 & -- & 4638 \\
huge\_sum\_512 & -- & 249.54 & -- & 4638 \\
\bottomrule
\end{tabular}
\end{center}


The benchmark uses seed 261008493, one shuffled warmup and five measured
rounds with identical controls. Its whole-cube ablation expands the entire
supplied braid, builds a reduced complex, produces a rank trace and replays
it. The hybrid uses the public source-bound forest interface, including its
transport checks, decomposition, factor proofs and replay. The ablation also
uses the host's transport and work-accounting conventions, with larger finite
allowances for its larger complex. Neither timed arm exhausts its allowances
on the compared cases. The explicit comparison is admitted only for inputs
with at most 12 crossings; larger inputs are omitted by preflight.

The two four-strand controls separately time the maintained scanner on an
already constructed PD diagram. That timing has no rank certificate, source
conversion or independent replay. It measures the practical cost of choosing
a straightforward proof-producing cube instead of the faster ordinary
recognizer; it is not a comparison between equivalent proof interfaces.
The mixed cases adjoin the trivial three-braid $ab$ through a singleton seam
to either a nine-crossing trivial or nontrivial four-braid. Larger mixed cases
replace $ab$ with its $u^{2^k}$ conjugate for $k=8,64,512$; the exceptional
factor remains nine crossings. No whole-input speedup ratio is inferred for
these enormous represented words. Raw samples and source hashes are in
\path{../fast/results/exceptional_cube_20261008.json}.

For the small trivial connected sum, whole-cube production and replay take
7.702 seconds at the median, versus 169.427 ms for the hybrid; the median
paired ratio is 44.278. For the nontrivial connected sum, the corresponding
values are 5.558 seconds, 139.172 ms and ratio 40.209. Both comparisons
include proof replay. The trivial-sum hybrid samples vary from 159 to 331 ms,
and its identical control has a different median despite a paired A/A ratio
near one. We therefore interpret these as roughly fortyfold family-specific
gains, not precise crossover thresholds. All variation is retained.

The large mixed examples require exactly 4,638 cumulative cube generators
for production and replay, independent of the compressed sleeve exponent.
At $k=512$ the complete hybrid median is 249.543 ms and its proof has
137,317 canonical bytes. The measured time grows with the compressed
source and proof, not its expanded conjugating length.

On the standalone nine-crossing four-braids, the hybrid cube path takes
161.835 ms for the unknot and 137.068 ms for the nontrivial knot. The ordinary
scanner takes 1.673 and 4.559 ms respectively, without independent proof
output. The cube is consequently a conservative complete fallback, not a
replacement for the faster scanner in ordinary diagram recognition.
These measurements motivated two further optimizations: verified explicit
reductions before a small factor's cube, and immediate termination on a proved
nontrivial factor. Section~\ref{sec:adaptive-forest} implements both with
source-bound proofs. Neither optimization is assumed in the checkpoint
measurements above.
